% RESUMO
A evasão no ensino superior raramente tem causa única, mas costuma ser acompanhada variável por variável. Este trabalho investiga se um grafo de semelhança entre estudantes ajuda a identificar perfis associados à evasão e ao sucesso, e qual combinação de distância e método faz isso melhor. Usamos o conjunto público UCI 697, com os 3.630 alunos que têm desfecho final (39,1\% de evasão). Cada aluno é um vértice, ligado aos mais parecidos por um grafo de $k$ vizinhos mútuos ($k = 40$); o desfecho nunca entra na construção. Quatro variações combinam duas distâncias para dados mistos (Gower e associação) com dois métodos (comunidades de Leiden e a rede neural em grafos GraphSAGE), em três janelas da trajetória: matrícula, fim do 1º semestre e fim do 2º semestre. A avaliação tem duas camadas: uma comum, em que as quatro variações são testadas nas mesmas 50 divisões 80/20 (F1 macro, AUC e revocação da evasão), e uma específica de cada método. O grafo é homofílico e as comunidades separam a evasão muito acima do acaso (V de Cramér de 0,32 a 0,62; $p = 0{,}001$). A GraphSAGE é a melhor variação em todas as janelas (F1 de 0,67, 0,86 e 0,89) e supera o Leiden em 98\% a 100\% das divisões; o Leiden, em troca, resume os alunos em poucos grupos interpretáveis. A distância por associação melhora as comunidades do Leiden, sobretudo no 1º semestre, e não altera a GraphSAGE. Na matrícula, as duas distâncias apontam o mesmo grupo de risco (alunos mais velhos, com pais de escolaridade desconhecida, cerca de 70\% de evasão), e o 1º semestre é o ponto de virada da trajetória.
